\section*{Resumen}
\addcontentsline{toc}{section}{Resumen}

Este protocolo de investigación doctoral aborda el estudio teórico y computacional de la ingeniería de interfaces y la modulación de propiedades topológicas en heteroestructuras compuestas por aleaciones de Heusler y materiales bidimensionales tipo MXene. Mediante simulaciones de primeros principios basadas en la Teoría de la Funcional de la Densidad (DFT), se investigará la estabilidad estructural, la hibridación electrónica en la interfaz y la inducción de estados topológicos no triviales (como aislantes topológicos bidimensionales o semimetales de Weyl/Dirac). El objetivo principal es sintonizar estas propiedades mediante efectos de deformación (strain), dopaje y la aplicación de campos eléctricos externos, proporcionando así una base fundamental para el diseño de dispositivos espintrónicos y topológicos de próxima generación.

\vspace{0.5cm}
\noindent \textbf{Palabras clave:} Heteroestructuras, Aleaciones de Heusler, MXenes, Ingeniería de interfaces, Modulación topológica, DFT.

\vspace{1.5cm}

\section*{Abstract}
\addcontentsline{toc}{section}{Abstract}

This doctoral research protocol focuses on the theoretical and computational study of interface engineering and the modulation of topological properties in heterostructures composed of Heusler alloys and two-dimensional MXene materials. Using first-principles simulations based on Density Functional Theory (DFT), we will investigate structural stability, electronic hybridization at the interface, and the induction of non-trivial topological states (such as two-dimensional topological insulators or Weyl/Dirac semimetals). The primary goal is to tune these properties through strain engineering, doping, and the application of external electric fields, thereby providing a fundamental basis for the design of next-generation spintronic and topological devices.

\vspace{0.5cm}
\noindent \textbf{Keywords:} Heterostructures, Heusler alloys, MXenes, Interface engineering, Topological modulation, DFT.
